\noindent\textbf{(a)} $e^y(y-1)=\sin x-x\cos x+C$, en sus arcos gráficos $C^1$.

\noindent\textbf{(b)} $y=-x-\log(C-x^2/2)$, con $C>0$ y $|x|<\sqrt{2C}$.

\noindent\textbf{(c)} $y=C(1+x)$; estas soluciones se extienden a todo $\mathbb R$.

\noindent\textbf{(d)} $y^2=x^2+Cx^4$, con $1+Cx^2>0$ para $x\ne0$. Las ramas pueden escribirse $y=\pm x\sqrt{1+Cx^2}$. En el origen se permiten constantes distintas a izquierda y derecha si coinciden las pendientes.

\noindent\textbf{(e)} $y=0$ o $y=x/(\log|x|+C)$ para $x\ne0$, con $\log|x|+C\ne0$. Las ramas tienen extensión $C^1$ al origen con $y(0)=y^{\prime}(0)=0$; pueden empalmarse allí con constantes independientes o con la rama nula.
